\section{Introduction}

A neural network that achieves 99\% predictive accuracy can fundamentally misunderstand the structure of its input---learning the right answers through the wrong mechanism. This counterintuitive finding lies at the heart of understanding when architectural inductive biases are essential versus merely convenient.

Weight-space learning---predicting model properties directly from weight tensors---has emerged as a promising paradigm for tasks ranging from accuracy prediction to model quality assessment \cite{schurholt2022hyper}. Central to this paradigm is the question of how to handle permutation symmetries: the weights of a neural network can be permuted across hidden neurons without changing the function it computes. Permutation-equivariant architectures like Neural Functional Networks (NFN) \cite{zhou2023neural} explicitly encode this symmetry, while standard MLPs must learn it from data.

The conventional wisdom holds that with sufficient data diversity, non-equivariant architectures should eventually learn to exploit these symmetries. After all, if the training distribution contains many permuted versions of similar weight configurations, an MLP should discover that permutation-related weights map to similar properties. This assumption has profound implications: if true, architectural inductive biases become a matter of convenience (improving sample efficiency) rather than necessity.

We challenge this assumption through systematic experiments across data scales from 1K to 40K models. Our key finding contradicts the ``data teaches invariance'' hypothesis: \textbf{MLPs trained on 40K models achieve near-perfect predictive accuracy ($R^2=0.99$) but fail to develop permutation invariance (0.63 vs NFN's 1.0)}. This dissociates task performance from mechanism---high accuracy can coexist with fundamentally incorrect representations that exploit dataset-specific position-accuracy correlations rather than semantic weight structure.

This insight reveals that permutation invariance must be built in architecturally; it cannot be learned from data regardless of scale. The implications extend beyond weight-space learning: certain symmetries may require explicit architectural encoding, with data quantity unable to substitute for structural inductive bias.

Our contributions are threefold:

\begin{enumerate}
    \item \textbf{Sample Efficiency of Equivariance:} We demonstrate that NFN achieves $R^2=0.95$ at N=1K while matched-capacity MLPs achieve only $R^2=0.35$---a 60 percentage point advantage that far exceeds expected margins, establishing the practical value of equivariant architectures in data-limited regimes.

    \item \textbf{Mechanism Verification:} We develop a probe invariance test that measures whether models develop permutation-invariant representations. Using this test, we show that NFN's invariance is mathematically perfect (deviation $< 1.19 \times 10^{-7}$ across permutations) while MLP invariance plateaus at 0.63 even at N=40K.

    \item \textbf{Falsification of Learned Invariance:} We provide the first direct evidence that MLPs cannot learn permutation invariance from data diversity. Despite achieving $R^2=0.986$ at large scale, MLPs learn position-sensitive statistics rather than permutation-invariant representations.
\end{enumerate}
